%% Conclusion (MedIA revision, v2 pipeline).
\section{Conclusion}
\label{sec:conclusion}

We asked whether pathology models can be trained once and then used at new
hospitals and scanners without retraining, and how a practitioner can tell.
Holding out every domain of four public cohorts, repeating every run and
selecting models without test data, we found that the train-once strategy can
work well: the best standard-scale method of each cohort kept a worst-site
AUROC between 0.88 and 0.99 without seeing labels from the test site. It is
not guaranteed, however. No method was best on every cohort, ERM fell to near
chance on one scanner, and with three to seven sites few differences between
methods could be confirmed, because most of the uncertainty comes from the
sites, not from the seeds.

Fine-tuned pathology foundation models with color augmentation were the most
reliable recipe on average, and color augmentation alone was close behind.
Re-estimating batch-normalization statistics on unlabeled images of the new
site helped mainly under severe scanner shift. Under the primary selection
rule, domain-generalization objectives never reached the top five, and frozen
foundation-model features were weak for mitotic figures. At the standard
scale, rankings agreed more between cohorts that share a task than between
cohorts that share a type of shift, so evidence from one task transfers poorly
to another.

Because development results cannot guarantee performance at a new site, we
proposed a site acceptance test that combines them with a small labeled sample
from that site. Validated on every held-out site, it kept false acceptance at
or below 2.1\% at standard scale and needed about half as many labels as local
labels alone, whereas the development results alone accepted 8.6\% of the
sites that fell below an AUROC of 0.80.

The roadmap of Section~\ref{sec:roadmap} turns these results into six steps:
define a worst-site target, collect many sites, train once with a strong
recipe, adapt without labels where the pipeline allows it, validate by holding
out whole sites, and accept each new site with the site acceptance test. The
protocol, the diagnostics and all per-run records are released so that other
cohorts and methods can be added to the same analysis.
